\documentclass{subfiles}
\begin{document}
    Throughout these notes we remind several concepts of both Algebra and Geometry.
        For the time being we focus on a brief Algebra reminder,
        in which we focus on some results on rings and ideals.

    \begin{definition*}[Rings]
        Let \(R\) be a set and 
        \[\begin{aligned}
            + &: R \times R \to R \\ 
            * &: R \times R \to R
        \end{aligned}\]
        be two operations. We say that \(R\) is a (commutative) ring if the following holds.
        \begin{itemize}
            \item For any \(a, b, c \in R, a + (b + c) = (a + b) + c\).
            \item Exists an element \(0 \in R\) such that \(a + 0 = 0 + a = a\).
            \item For any \(a \in R\), there is \(-a \in R\) such that \(a + (-a) = -a + a = 0\).
            \item For any \(a, b \in R, a + b = b + a\).
            \item For any \(a, b, c \in R, a * (b * c) = (a * b) * c\).
            \item For all \(a, b \in R, a * b = b * a\).
            \item \(*\) distributes over \(+\).
        \end{itemize}
    \end{definition*}

    \begin{definition*}[Ideal]
        Let \(R\) be a ring. We say that a subset \(I \subseteq R\) is and ideal 
            of \(R\) if it holds that 
            \begin{itemize}
                \item \(I\) is subgroup for \((R, +)\). That is, if 
                    \[
                        \forall x, y \in I, x + y \in I.
                    \]
                \item \(\forall a \in R, \forall x \in I, a * x \in I\).
            \end{itemize}
    \end{definition*}

    \begin{remark*}
        From now on we write \(ab\) instead of \(a * b\).
    \end{remark*}

    \begin{definition*}
        Let \(R\) be a ring and let \(T \subseteq R\) be a subset. 
            We define the set 
            \[
                (T) = \set*{\sum_{i = 0}^{n} a_{i}x_{i}}[a_{i} \in R, x_{i} \in T].
            \]

        Additionally, if \(I = (T)\) and \(\abs{T} < \infty\) we say that \(I\)
            is \emph{finetely generated}.
    \end{definition*}

    \begin{definition*}
        Let \(R, Q\) be rings. A function \(f : R \to Q\) is called (ring) homomorphism
            if 
            \begin{itemize}
                \item \(f(a + b) = f(a) + f(b), \forall a, b \in R\); \\ 
                \item \(f(ab) = f(a)f(b), \forall a, b \in R\); \\ 
                \item \(f(1_{R}) = 1_{Q}\).
            \end{itemize}
    \end{definition*}

    \begin{remark*}
        Given \(f : R \to Q\) a ring homomorphism, it can be proved that 
            \[
                \ker{f} = \set*{a \in R}[f(a) = 0]
            \]
            is an ideal of \(R\). On the other hand, 
            \(\im{f} = \set{b \in Q}[b = f(a), a \in R]\) is a subring of \(Q\).
    \end{remark*}

    Recall that, given \(R\) a ring a \(I \subseteq R\) an ideal, 
        we can consider the quotient ring 
        \[
            \frac{A}{I} = \set*{a + I}[a \in R].
        \]
        We usually write \(\overline{a}\) instead of \(a + I\).

    Lastly, we present a classic result in ring theory,
    \begin{theorem*}[First isomorphism theorem]
        Given a homomorphism \(f : R \to Q\), there exists a unique homomorphism \(g\)
        such that \(f = \pi \circ g\), where \(\pi\) is the canonical epimorphism.
        In terms of commutative diagrams, we have the following.
        \subfile{../../extra/TikZ/Figure - Commutative diagram of FIT}
    \end{theorem*}

    \begin{proposition*}
        Given \(R\) a ring and \(\mathfrak{p}\) an ideal, we say that \(\mathfrak{p}\)
        is a prime ideal of \(R\), if 
        \[
            ab \in \mathfrak{p} \implies a \in \mathfrak{p} \lor b \in \mathfrak{p}
        \]
        or equivalently, if \(a, b \notin \mathfrak{p} \implies ab \notin \mathfrak{p}\).
    \end{proposition*}
\end{document}
